\chapter{Circuits and Proof Complexity}\label{ch:circuits}

\begin{center}\itshape
Almost every function is hard, and we cannot name one.\\
Here is how far the hunt has gone, and what stands in its way.\\
\upshape(the content of this chapter, in our words)
\end{center}

In 1937 Claude Shannon, a twenty-one-year-old student at the Massachusetts Institute of Technology, noticed that the relay circuits of telephone exchanges obey the rules of Boolean algebra, which he had met in a philosophy course; his master's thesis, published in 1938, turned the design of switching circuits into algebra. Twelve years later, in 1949, he asked how many relays a circuit needs, and answered with a counting argument of a few lines: there are many more Boolean functions than small circuits, so almost every function of $n$ variables needs a circuit of size about $2^n/n$. Then came the embarrassment that still lasts: nobody could exhibit a single explicit function that needs even a circuit of size $n^2$.

Circuits looked like the right battlefield for $\mathsf P$ versus $\mathsf{NP}$. A polynomial-time machine gives polynomial-size circuits, so to separate the classes it would suffice to show that some problem in $\mathsf{NP}$ needs large circuits; and circuits are finite, static, combinatorial objects, with no Turing machine and no diagonalisation in sight. In the early 1980s Michael Sipser and others set out on this programme, and the first results came fast. In 1981 Merrick Furst, James Saxe and Sipser, and independently Miklós Ajtai, proved that the parity function cannot be computed by circuits of constant depth and polynomial size. In 1985 Alexander Razborov, a young mathematician from Moscow, proved that detecting a clique needs monotone circuits of superpolynomial size, a result that Noga Alon and Ravi Boppana soon sharpened to an exponential bound; in 1987 Razborov and Roman Smolensky extended the parity bound to circuits with counting gates. Razborov received the Nevanlinna Prize in 1990.

Then the programme stalled, and in 1994 Razborov himself, with Steven Rudich, explained why: almost all the known lower bound proofs are \emph{natural}, and natural proofs cannot separate $\mathsf P$ from $\mathsf{NP}$ if cryptography is possible. The same years saw the birth of a twin programme, \emph{proof complexity}: in 1979 Stephen Cook and Robert Reckhow showed that $\mathsf{NP}\ne\mathsf{coNP}$ is equivalent to the statement that no proof system has short proofs of all tautologies, and in 1985 Armin Haken proved that resolution, the proof system of Tome~C1, Chapter~1, needs exponentially long proofs of the pigeonhole principle. This last chapter tells both stories.

\begin{think}
Almost every Boolean function of $100$ variables needs a circuit with more than $10^{27}$ gates. Can you name one? Can you name one that needs more than $10^4$ gates?
\end{think}

\section{Circuits}\label{sec:circuits}

\begin{definition}\label{def:circuit}
A \emph{Boolean circuit} with inputs $x_1,\dots,x_n$ is a directed acyclic graph whose sources are the inputs and the constants $0,1$, and whose other vertices, the \emph{gates}, are labelled $\land$ or $\lor$ (with two incoming arcs) or $\lnot$ (with one); one gate is the \emph{output}. Each gate computes the function obtained by applying its label to its predecessors. The \emph{size} of the circuit is its number of gates, its \emph{depth} the length of a longest path. A language $L\subseteq\{0,1\}^*$ is decided by a family $(C_n)$ of circuits if $C_n$ has $n$ inputs and computes the indicator of $L\cap\{0,1\}^n$.
\end{definition}

A formula is a circuit whose gates have one outgoing arc: circuits may reuse intermediate results, formulas may not. Figure~\ref{fig:circuits} shows two small circuits; the majority circuit on the left has the smallest possible size among circuits with gates computing arbitrary functions of two bits (the companion notebook checks it by exhaustive search).

\begin{figure}[ht]
\centering
\fig{ch12-circuits}
\caption{Left: a circuit of size $4$ for the majority of three bits, $(x_1\land x_2)\lor\bigl(x_3\land(x_1\lor x_2)\bigr)$. Right: parity of three bits with two exclusive-or gates (each of which costs a few gates of $\land,\lor,\lnot$). Every CNF for the parity of $n$ bits has at least $2^{n-1}$ clauses (Tome~C1, Proposition~1.5), but its circuits have linear size.}
\label{fig:circuits}
\end{figure}

\begin{proposition}\label{prop:ppoly}
A language is in $\mathsf P/\mathsf{poly}$ (Chapter~\ref{ch:randomness}) if and only if it is decided by a family of circuits of polynomial size. In particular $\mathsf P\subseteq\mathsf P/\mathsf{poly}$.
\end{proposition}

\begin{proof}
Suppose that $L\in\mathsf P/\mathsf{poly}$, with a machine $V$ running in time $t(n)$, polynomial, on inputs $\langle x,a_{|x|}\rangle$. In the proof of Theorem~\ref{thm:cooklevin} the computation of $V$ is a tableau of $t\times t$ cells, and the content of a cell at step $k+1$ is determined by the contents of the three cells above it at step $k$. Code each cell content by a constant number of bits; each bit of a cell at step $k+1$ is then a fixed Boolean function of a constant number of bits at step $k$, computed by a circuit of constant size (Tome~C1, Theorem~1.3). Stacking these circuits row after row, with the first row fed by $x$ and by the bits of $a_n$ wired as constants, gives a circuit of size $O(t^2)$ deciding $L\cap\{0,1\}^n$. Conversely, if $C_n$ decides $L\cap\{0,1\}^n$ and has polynomial size, take $a_n$ to be the description of $C_n$ and let $V$ evaluate the circuit $a_{|x|}$ on $x$, gate after gate, in polynomial time.
\end{proof}

So $\mathsf P\ne\mathsf{NP}$ would follow from a superpolynomial circuit lower bound for a single problem of $\mathsf{NP}$. Lower bounds exist in abundance; the trouble is that they are not explicit.

\begin{theorem}[Shannon 1949]\label{thm:shannon}
Let $n\ge7$ and $s=2^n/(2n)$. The fraction of Boolean functions of $n$ variables computed by some circuit of size at most $s$ is at most $2^{\,n-s(2\log_2n-\log_23)}$, which tends to $0$ very fast.
\end{theorem}

\begin{proof}
A circuit with at most $s$ gates can be padded with useless gates to exactly $\lfloor s\rfloor$ gates. Describe it by listing, for each gate, its label (three choices) and its at most two predecessors, among the $n$ inputs, the $2$ constants and the other gates; then name the output gate. Since $2^n\ge2n^2+4n$ for $n\ge7$, we have $n+2\le s$, so there are at most $n+2+s\le2s$ choices per predecessor, and the number of functions computed by circuits of size at most $s$ is at most
\[
s\,\bigl(3(2s)^2\bigr)^{s}=s\cdot2^{\,s(\log_212+2\log_2s)}.
\]
As $\log_2s=n-\log_2(2n)=n-1-\log_2n$, the exponent is $s(\log_212-2+2n-2\log_2n)=2^n+s(\log_23-2\log_2n)$, using $2ns=2^n$. Dividing by the number $2^{2^n}$ of functions, and bounding $s\le2^n$, gives the claim. For $n\ge2$ we have $2\log_2n-\log_23>0.4$, so the fraction is at most $2^{n-0.4s}$.
\end{proof}

The bound is sharp: Oleg Lupanov showed in 1958 that every function of $n$ variables has a circuit of size $(1+o(1))2^n/n$ (quoted); a simpler construction of size $O(2^n)$ is Exercise~\ref{ex:everyfn}. And yet the best lower bound known for any explicit function, one computable in polynomial time, is about $3.1n$ gates, obtained by Jiatu Li and Tianqi Yang in 2022 after decades of slow progress from $2n$ to $3n$ (quoted). Counting proves existence and gives no names.

\medskip
Would small circuits for $\mathsf{NP}$ be absurd? Not obviously: circuits can be non-uniform, with a different clever circuit for each length. But Richard Karp and Richard Lipton showed in 1980 that they would make the polynomial hierarchy collapse.

\begin{theorem}[Karp, Lipton 1980]\label{thm:karplipton}
If $\mathsf{NP}\subseteq\mathsf P/\mathsf{poly}$, then $\mathsf{PH}=\Sigma_2^p$.
\end{theorem}

\begin{proof}
By Theorem~\ref{thm:collapse} it is enough to show $\Pi_2^p\subseteq\Sigma_2^p$; then, taking complements, $\Sigma_2^p\subseteq\Pi_2^p$ and the two classes are equal. Let $L\in\Pi_2^p$:
\[
x\in L\iff\forall y\ \exists z\ V(x,y,z),
\]
with $y,z$ of polynomial length and $V$ polynomial. By the Cook--Levin reduction (Theorem~\ref{thm:cooklevin}), from $(x,y)$ one computes in polynomial time a formula $\varphi_{x,y}$ that is satisfiable if and only if $\exists z\,V(x,y,z)$; pad all these formulas, for a given $|x|$, to a common length $m$, polynomial in $|x|$, in a coding where replacing a variable by a constant keeps the length.

Since \textup{\textsc{Sat}} is in $\mathsf P/\mathsf{poly}$, there are circuits $D_m$ of polynomial size that decide the satisfiability of formulas of length $m$. They can be turned into circuits $E_m$ that \emph{find} satisfying assignments, by self-reducibility (Proposition~\ref{prop:selfreduce}): given $\varphi$ with variables $u_1,\dots,u_v$, substitute $u_1=1$, ask $D_m$ whether the result is satisfiable, keep $u_1=1$ if so and set $u_1=0$ otherwise, and continue with $u_2$. Each step is a copy of $D_m$ and a substitution circuit of polynomial size, so $E_m$ has polynomial size $r(|x|)$, and $E_m(\varphi)$ satisfies $\varphi$ whenever $\varphi$ is satisfiable. Now
\[
x\in L\iff\exists E\ \text{of size at most }r(|x|)\ \ \forall y\ \ E(\varphi_{x,y})\text{ satisfies }\varphi_{x,y}.
\]
If $x\in L$, the circuit $E_m$ works. Conversely, if some circuit $E$ produces, for every $y$, a satisfying assignment of $\varphi_{x,y}$, then all these formulas are satisfiable and $x\in L$. The inner condition is checked in polynomial time, so $L\in\Sigma_2^p$.
\end{proof}

So circuit lower bounds are expected: the collapse of the hierarchy is considered almost as unlikely as $\mathsf P=\mathsf{NP}$. The next three sections prove lower bounds for restricted circuits, where the hunt succeeded.

\section{Constant depth: parity}\label{sec:parity}

Allow gates $\land$ and $\lor$ of unbounded fan-in, but bound the depth. The class $\mathsf{AC}^0$ consists of the languages decided by circuits of constant depth and polynomial size, built from unbounded fan-in $\land,\lor$ and $\lnot$; adding gates $\mathrm{MOD}_3$, which output $1$ when the number of their inputs equal to $1$ is \emph{not} divisible by $3$, gives the class $\mathsf{AC}^0[3]$. Addition of two numbers is in $\mathsf{AC}^0$ (by a carry-lookahead formula), and so is the recognition of every regular language that has no counting modulo; parity is the archetype of what counting needs. We write $\mathrm{PAR}_n(x)=x_1\oplus\dots\oplus x_n$.

The proof of Razborov and Smolensky approximates circuits by polynomials over the field $\mathbb F_3=\{0,1,2\}$, in two steps: small circuits are close to low-degree polynomials, and parity is far from them. Inputs $x\in\{0,1\}^n$ are seen in $\mathbb F_3^n$.

\begin{lemma}[small circuits are almost polynomials]\label{lem:approxpoly}
Let $C$ be a circuit of depth $d$ and size $s$ with gates $\land,\lor,\lnot,\mathrm{MOD}_3$ of unbounded fan-in, and let $\ell\ge1$. There is a polynomial $p\in\mathbb F_3[x_1,\dots,x_n]$ of degree at most $(2\ell)^d$ such that $p(x)=C(x)$ for all but at most $s2^{-\ell}2^n$ points $x\in\{0,1\}^n$.
\end{lemma}

\begin{proof}
We attach to each gate $g$ a \emph{random} polynomial $p_g$, by induction along the circuit, starting from $p=x_i$ at the input $x_i$, and show two things: $\deg p_g\le(2\ell)^{\mathrm{depth}(g)}$, and for each fixed $x$, if the polynomials of the predecessors of $g$ are correct at $x$ (equal to the value of the gate), then $p_g$ is correct at $x$ with probability at least $1-2^{-\ell}$, and with probability $1$ for gates $\lnot$ and $\mathrm{MOD}_3$.

For $g=\lnot h$, take $1-p_h$. For $g=\mathrm{MOD}_3(h_1,\dots,h_k)$, take $\bigl(\sum_ip_{h_i}\bigr)^2$: in $\mathbb F_3$ the square of $0$ is $0$ and the squares of $1$ and $2$ are $1$. The degree doubles. For $g=\lor(h_1,\dots,h_k)$, choose independent uniformly random subsets $S_1,\dots,S_\ell$ of $\{1,\dots,k\}$ and take
\[
p_g=1-\prod_{j=1}^{\ell}\Bigl(1-\Bigl(\sum_{i\in S_j}p_{h_i}\Bigr)^2\Bigr),
\]
of degree at most $2\ell$ times the maximum degree of the $p_{h_i}$. Fix $x$ and suppose the $p_{h_i}$ are correct at $x$. If all the $h_i(x)$ are $0$, all the sums vanish and $p_g(x)=0$. If some $h_{i_0}(x)=1$, pair each subset $S$ with $S\triangle\{i_0\}$: the two sums $\sum_{i\in S}h_i(x)$ differ by $1$, so at most one of them is $0$ in $\mathbb F_3$, and a random $S_j$ gives a sum $0$ with probability at most $\frac12$. Unless all $\ell$ sums vanish, an event of probability at most $2^{-\ell}$, some factor is $0$ and $p_g(x)=1$. Gates $\land$ are treated through $\land(h_i)=\lnot\lor(\lnot h_i)$, with the same degree and error.

Each gate multiplies the degree by at most $2\ell$, so $\deg p\le(2\ell)^d$ for the output polynomial $p$. For a fixed $x$, the output is wrong only if some gate errs at $x$ while its predecessors are right, so by the union bound $P\bigl(p(x)\ne C(x)\bigr)\le s2^{-\ell}$. The expected number of bad points is at most $s2^{-\ell}2^n$, so some choice of the random subsets does at least as well.
\end{proof}

\begin{lemma}[parity is far from low degree]\label{lem:parityfar}
Let $p\in\mathbb F_3[x_1,\dots,x_n]$ have degree at most $D$. Then $p(x)=\mathrm{PAR}_n(x)$ for at most $\sum_{i\le n/2+D}\binom ni$ points $x\in\{0,1\}^n$, which is at most $2^n\bigl(\frac12+\frac{D+1}{\sqrt n}\bigr)$.
\end{lemma}

\begin{proof}
Change variables: $y_i=1+x_i$ maps $\{0,1\}$ to $\{1,-1\}$ in $\mathbb F_3$ (as $2=-1$), and then $\prod_iy_i=(-1)^{x_1+\dots+x_n}=1+\mathrm{PAR}_n(x)$. So $q(y)=1+p(y_1-1,\dots,y_n-1)$, of degree at most $D$, equals $\prod_iy_i$ on the set $G'\subseteq\{\pm1\}^n$ that corresponds to the set $G$ of points where $p$ agrees with parity.

Every function $f:G'\to\mathbb F_3$ is the restriction of a polynomial, for instance $\sum_af(a)\prod_i\bigl(-(1+a_iy_i)\bigr)$, since $-(1+a_iy_i)$ is $1$ when $y_i=a_i$ and $0$ when $y_i=-a_i$ (recall $-2=1$ in $\mathbb F_3$). As $y_i^2=1$ on $\{\pm1\}$, we can make it multilinear, a combination of monomials $y_I=\prod_{i\in I}y_i$. If $|I|>n/2$, then on $\{\pm1\}^n$ we have $y_I=y_{\{1,\dots,n\}}\,y_{\overline I}$, where $\overline I$ is the complement, and on $G'$ this is $q(y)\,y_{\overline I}$, of degree at most $D+n-|I|<D+n/2$. After making it multilinear again, $f$ is on $G'$ a multilinear polynomial of degree at most $n/2+D$. Different functions on $G'$ need different polynomials, so
\[
3^{|G|}\le3^{\#\{I:\ |I|\le n/2+D\}},\qquad |G|\le\sum_{i\le n/2+D}\binom ni.
\]
For the estimate, the terms $i<n/2$ sum to at most $2^{n-1}$, by the symmetry $\binom ni=\binom n{n-i}$, and there are at most $D+1$ other terms, each at most $\binom n{\lfloor n/2\rfloor}\le2^n/\sqrt n$. This last inequality follows from $\binom{2m}m\le4^m/\sqrt{3m+1}$, proved by induction since $\bigl(\frac{2m+1}{2m+2}\bigr)^2\le\frac{3m+1}{3m+4}$ (expand: it says $12m^3+28m^2+19m+4\le12m^3+28m^2+20m+4$), and from $\binom{2m+1}m=\frac12\binom{2m+2}{m+1}$.
\end{proof}

\begin{theorem}[Razborov 1987, Smolensky 1987]\label{thm:rs}
For $n\ge400$, every circuit of depth $d$ with gates $\land,\lor,\lnot,\mathrm{MOD}_3$ of unbounded fan-in that computes $\mathrm{PAR}_n$ has size at least $\frac15\,2^{\,n^{1/(2d)}/40}$. In particular parity is not in $\mathsf{AC}^0[3]$, nor in $\mathsf{AC}^0$.
\end{theorem}

\begin{proof}
Let $\ell$ be the largest integer with $(2\ell)^d\le\sqrt n/10-1$; if $\ell=0$, the bound to prove is at most $\frac25<1$ and there is nothing to do. Lemma~\ref{lem:approxpoly} gives a polynomial of degree $D\le\sqrt n/10-1$ agreeing with the circuit, hence with parity, on at least $(1-s2^{-\ell})2^n$ points, and Lemma~\ref{lem:parityfar} says that it agrees on at most $(\frac12+\frac1{10})2^n$. So $s2^{-\ell}\ge\frac25$, that is $s\ge\frac25\,2^\ell$. Finally $\sqrt n/10-1\ge\sqrt n/20$ for $n\ge400$, so $2\ell+2>(\sqrt n/20)^{1/d}\ge n^{1/(2d)}/20$, whence $\ell\ge n^{1/(2d)}/40-1$ and $s\ge\frac15\,2^{n^{1/(2d)}/40}$.
\end{proof}

Nothing in the proof is special to parity and $3$: the same argument shows that $\mathrm{MOD}_q$ is not in $\mathsf{AC}^0[p]$ for distinct primes $p,q$ (quoted). Johan Håstad's \emph{switching lemma} of 1986 gives the sharp bound $2^{\Omega(n^{1/(d-1)})}$ for $\mathsf{AC}^0$ circuits of depth $d$ (quoted): a random restriction, fixing most variables at random, turns a small CNF into a short decision tree with high probability, so the bottom two layers of the circuit can be switched and merged, and depth decreases while parity of the unfixed variables survives. For gates modulo composite numbers, such as $\mathrm{MOD}_6$, the polynomial method breaks down, and the class $\mathsf{ACC}^0$ remained mysterious until Ryan Williams proved in 2011 that it does not contain $\mathsf{NEXP}$ (quoted), with an unexpected weapon: a satisfiability algorithm for $\mathsf{ACC}^0$ circuits slightly faster than exhaustive search.

\section{Monotone circuits: cliques}\label{sec:monotone}

A circuit is \emph{monotone} if it has no $\lnot$ gates. Monotone circuits compute exactly the monotone functions (Tome~C1, Proposition~1.4 and Exercise~1.2), such as
\[
\mathrm{CLIQUE}_{k,n}(G)=1\iff\text{the graph }G\text{ contains a clique of size }k,
\]
whose $\binom n2$ inputs are the indicators $x_{uv}$ of the edges of a graph $G$ on the vertex set $\{1,\dots,n\}$. For each $k$ this function is $\mathsf{NP}$-complete when $k$ is part of the input (Theorem~\ref{thm:isclique}).

Razborov's \emph{approximation method} replaces each gate of a small monotone circuit by a crude approximation from a small family, and shows that the approximations make few errors on two families of test graphs, whereas every member of the family makes many. The tool for keeping the approximations crude is the sunflower lemma of Erd\H os and Rado (Tome~C1, Theorem~8.18): a family of more than $\ell!(p-1)^\ell$ sets of size at most $\ell$ contains a \emph{sunflower} with $p$ petals, that is $p$ sets $Z_1,\dots,Z_p$ whose pairwise intersections all equal the same \emph{core} $Z$.

\emph{The test graphs.} A \emph{positive} test graph is a clique on a uniformly random set $K$ of $k$ vertices, with no other edge; $\mathrm{CLIQUE}_{k,n}$ is $1$ on it. A \emph{negative} test graph comes from a uniformly random colouring $c:\{1,\dots,n\}\to\{1,\dots,k-1\}$: two vertices are joined when they have different colours. By the pigeonhole principle it has no $k$-clique, and $\mathrm{CLIQUE}_{k,n}$ is $0$ on it.

\emph{The approximators.} For $S\subseteq\{1,\dots,n\}$ let $C_S(G)=1$ if $S$ is a clique of $G$ (so $C_S=1$ when $|S|\le1$). Fix $\ell,p$ and let $M=(p-1)^\ell\ell!$. An \emph{approximator} is a disjunction $C_{S_1}\lor\dots\lor C_{S_r}$ with $r\le M$ and $|S_i|\le\ell$ (for $r=0$ it is the constant $0$). On a negative test graph, $C_S=1$ means that $S$ is \emph{rainbow}: its vertices have distinct colours.

\begin{theorem}[Razborov 1985; Alon, Boppana 1987]\label{thm:monotoneclique}
Let $\ell\ge2$ and $k=\ell^2\le n^{1/4}$. For $n$ large enough, every monotone circuit computing $\mathrm{CLIQUE}_{k,n}$ has at least $n^{\ell/5}$ gates. With $\ell=\lfloor n^{1/8}\rfloor$, this is $2^{\Omega(n^{1/8}\log n)}$.
\end{theorem}

\begin{proof}
Take $p=\lceil10\ell\log_2n\rceil$. We build, gate by gate, an approximator $A_g$ for each gate $g$ of a monotone circuit of size $s$. At an input $x_{uv}$, $A=C_{\{u,v\}}$, exactly. To \emph{pluck} a family of more than $M$ sets, find a sunflower $Z_1,\dots,Z_p$ in it, with core $Z$, and replace its petals by $Z$ (Figure~\ref{fig:sunflower}); repeat until at most $M$ sets remain. Each plucking removes at least $p-1\ge1$ sets.
\begin{itemize}
\item $g=h\lor h'$: take the union of the two families of $A_h$ and $A_{h'}$, and pluck.
\item $g=h\land h'$: if $A_h=\bigvee_iC_{S_i}$ and $A_{h'}=\bigvee_jC_{T_j}$, take the family of all $S_i\cup T_j$, \emph{drop} those of size more than $\ell$, and pluck.
\end{itemize}
Before plucking there are at most $M^2$ sets, so a gate performs at most $M^2$ pluckings and drops at most $M^2$ sets.

\emph{Positive test graphs.} Let $Y$ be a clique on $K$. The step $C_{S_i}\land C_{T_j}\to C_{S_i\cup T_j}$ is exact on $Y$: both sides say $S_i\cup T_j\subseteq K$. Plucking only increases the approximator, since $Z\subseteq Z_i$ gives $C_Z\ge C_{Z_i}$. So by induction along the circuit $A_g(Y)\ge g(Y)$ for every gate, unless at some gate a dropped set $R$, of size at least $\ell+1$, satisfies $R\subseteq K$. For a fixed $R$, this has probability
\[
\frac{\binom{n-|R|}{k-|R|}}{\binom nk}=\frac{k(k-1)\cdots(k-|R|+1)}{n(n-1)\cdots(n-|R|+1)}\le\Bigl(\frac kn\Bigr)^{\ell+1}.
\]
So $A_{\mathrm{out}}(Y)=1$ for all positive test graphs except a fraction at most $sM^2(k/n)^{\ell+1}$.

\emph{Negative test graphs.} Let $N$ come from a colouring $c$. The step $C_{S_i}\land C_{T_j}\to C_{S_i\cup T_j}$ and the dropping can only decrease the value. A plucking can increase it from $0$ to $1$ only if $C_Z(N)=1$ while $C_{Z_1}(N)=\dots=C_{Z_p}(N)=0$. Fix the colours of $Z$, rainbow. The sets $Z_i\setminus Z$ are disjoint, so the events ``$Z_i$ is not rainbow'' are independent; each requires one of the at most $\binom\ell2$ pairs of $Z_i$ not inside $Z$ to have equal colours, and a pair with a vertex outside $Z$ has equal colours with probability $\frac1{k-1}$. As $\binom\ell2=\frac{k-\ell}2\le\frac{k-1}2$, each event has probability at most $\frac12$, and all of them together at most $2^{-p}$. So by induction $A_g(N)\le g(N)$ for every gate, unless some plucking errs at $N$, and $A_{\mathrm{out}}(N)=0$ for all negative test graphs except a fraction at most $sM^2 2^{-p}$.

\emph{Every approximator errs a lot.} If $A_{\mathrm{out}}$ is the constant $0$, it is wrong on all positive test graphs, so $sM^2(k/n)^{\ell+1}\ge1$. Otherwise it contains some $C_S$ with $|S|\le\ell$, and $S$ is rainbow with probability at least $1-\binom\ell2/(k-1)\ge\frac12$, so $A_{\mathrm{out}}=1$ on at least half of the negative test graphs, and $sM^22^{-p}\ge\frac12$. In both cases
\[
s\ge\frac{\min\bigl((n/k)^{\ell+1},\,2^{p-1}\bigr)}{M^2}.
\]
Now $M\le(p\ell)^\ell$ and $p\ell\le11\ell^2\log_2n=11k\log_2n\le n^{1/4+o(1)}$, so $M^2\le n^{(1/2+o(1))\ell}$; while $(n/k)^{\ell+1}\ge n^{3\ell/4}$ and $2^{p-1}\ge\frac12n^{10\ell}$. Hence $s\ge n^{(1/4-o(1))\ell}\ge n^{\ell/5}$ for $n$ large.
\end{proof}

\begin{figure}[ht]
\centering
\fig{ch12-sunflower}
\caption{Plucking. The petals $Z_1,\dots,Z_p$ of a sunflower meet pairwise in the core $Z$. Replacing $C_{Z_1}\lor\dots\lor C_{Z_p}$ by $C_Z$ shrinks the family; it never hurts on positive test graphs, and on a negative test graph it errs only if the core is rainbow while all $p$ petals, coloured independently outside the core, fail to be: probability at most $2^{-p}$.}
\label{fig:sunflower}
\end{figure}

Razborov's original bound was $n^{\Omega(\log n)}$; Alon and Boppana's choice of parameters, which we followed, gives an exponential one. Why does this not separate $\mathsf P$ from $\mathsf{NP}$? Because negations matter: Razborov proved the same year that the perfect matching function, decidable in polynomial time (Tome~C1, Chapter~6), also needs monotone circuits of size $n^{\Omega(\log n)}$, and \'Eva Tardos showed in 1988 that a monotone function computable in polynomial time can need exponential monotone circuits (quoted). Monotone lower bounds say nothing about general circuits.

\section{Proof complexity}\label{sec:proofcomplexity}

A formula is unsatisfiable when its negation is a tautology. How long must a proof of unsatisfiability be?

\begin{definition}\label{def:refsystem}
A \emph{refutation system} is a polynomial-time algorithm $V$ such that a CNF formula $\varphi$ is unsatisfiable if and only if $V(\varphi,\pi)$ accepts for some string $\pi$, a \emph{refutation} of $\varphi$. It is \emph{polynomially bounded} if every unsatisfiable $\varphi$ has a refutation of length polynomial in $|\varphi|$.
\end{definition}

Resolution is a refutation system, by its soundness and completeness (Tome~C1, Definition~1.10 and Theorem~1.12): $\pi$ is a list of clauses ending with the empty clause, each one in $\varphi$ or a resolvent of two earlier ones.

\begin{proposition}[Cook, Reckhow 1979]\label{prop:cookreckhow}
There is a polynomially bounded refutation system if and only if $\mathsf{NP}=\mathsf{coNP}$.
\end{proposition}

\begin{proof}
The set of unsatisfiable CNF formulas is $\mathsf{coNP}$-complete, as the complement of the $\mathsf{NP}$-complete \textup{\textsc{Sat}} (Theorem~\ref{thm:cooklevin}; a formula is first made a CNF by the Tseitin transformation, Tome~C1, Theorem~1.6). A polynomially bounded refutation system is exactly a polynomial verifier with short certificates for this set, so it exists if and only if the set is in $\mathsf{NP}$, that is if and only if $\mathsf{coNP}\subseteq\mathsf{NP}$, which is equivalent to $\mathsf{NP}=\mathsf{coNP}$.
\end{proof}

So proving superpolynomial lower bounds for stronger and stronger refutation systems is a programme towards $\mathsf{NP}\ne\mathsf{coNP}$, hence towards $\mathsf P\ne\mathsf{NP}$. Its first great success concerns the pigeonhole formula $\mathrm{PHP}_n$ of Tome~C1, Exercise~1.5: variables $p_{ij}$ (pigeon $i$ sits in hole $j$) for $n+1$ pigeons and $n$ holes, \emph{pigeon clauses} $p_{i1}\lor\dots\lor p_{in}$, and \emph{hole clauses} $\lnot p_{ij}\lor\lnot p_{i'j}$ for $i\ne i'$.

The proof below, by Paul Beame and Toniann Pitassi (1996), studies refutations through \emph{critical assignments}: an assignment is \emph{$i$-critical} if it puts the $n$ pigeons other than $i$ into the $n$ holes bijectively, and pigeon $i$ nowhere. It satisfies every clause of $\mathrm{PHP}_n$ except the pigeon clause of $i$. On critical assignments a negative literal can be made positive: $\lnot p_{ij}$ holds if and only if hole $j$ is occupied by another pigeon. So for a clause $C$ let $C^+$ be the clause obtained by replacing each literal $\lnot p_{ij}$ by $\bigvee_{i'\ne i}p_{i'j}$; it is \emph{monotone} (a disjunction of variables), and it agrees with $C$ on every critical assignment. Write $|C^+|$ for its number of variables.

\begin{lemma}[Beame, Pitassi]\label{lem:bp}
Let $m\ge4$, and let $C_1,\dots,C_S$ be clauses in the variables of $\mathrm{PHP}_m$ such that each $C_t$ is implied by a clause of $\mathrm{PHP}_m$ or by two earlier clauses, and $C_S$ is empty. Then some $C_t$ satisfies $|C_t^+|\ge m^2/9$.
\end{lemma}

\begin{proof}
Let $\mathrm{pig}(C)$ be the set of pigeons $i$ such that some $i$-critical assignment falsifies $C$. A clause implied by a clause of $\mathrm{PHP}_m$ has $|\mathrm{pig}|\le1$; the empty clause has all $m+1$ pigeons; and if $C_t$ is implied by $C_u$ and $C_v$, an assignment falsifying $C_t$ falsifies $C_u$ or $C_v$, so $\mathrm{pig}(C_t)\subseteq\mathrm{pig}(C_u)\cup\mathrm{pig}(C_v)$. Let $C$ be the first clause with $|\mathrm{pig}(C)|\ge m/3$. It is not implied by a clause of $\mathrm{PHP}_m$ (as $m/3>1$), so its two premises have fewer than $m/3$ pigeons each, and $m/3\le|\mathrm{pig}(C)|<2m/3$.

Let $i\in\mathrm{pig}(C)$ and $j\notin\mathrm{pig}(C)$. Take an $i$-critical assignment $\alpha$ falsifying $C$, and let $l$ be the hole of pigeon $j$ in $\alpha$. Move $j$ out and $i$ into hole $l$: the result $\alpha'$ is $j$-critical (Figure~\ref{fig:pigeons}), so it satisfies $C$, as $j\notin\mathrm{pig}(C)$. Thus $C^+$ is $0$ at $\alpha$ and $1$ at $\alpha'$. The two assignments differ only in $p_{il}$, which goes from $0$ to $1$, and $p_{jl}$, which goes from $1$ to $0$; since $C^+$ is a disjunction of variables, it must contain $p_{il}$. For a fixed $i$, different pigeons $j\notin\mathrm{pig}(C)$ give different holes $l$, since $\alpha$ is injective. So $C^+$ contains at least $|\mathrm{pig}(C)|\cdot(m+1-|\mathrm{pig}(C)|)\ge\frac m3\cdot\frac m3$ variables.
\end{proof}

\begin{figure}[ht]
\centering
\fig{ch12-pigeons}
\caption{The swap in the proof of Lemma~\ref{lem:bp}, for $5$ holes and $6$ pigeons (rows are pigeons, columns holes, filled cells the variables equal to $1$). Moving pigeon $j$ out of hole $l$ and pigeon $i$ in changes only the two framed variables; a monotone clause that is false before and true after must contain $p_{il}$.}
\label{fig:pigeons}
\end{figure}

\begin{theorem}[Haken 1985]\label{thm:haken}
For $n\ge13$, every resolution refutation of $\mathrm{PHP}_n$ has at least $(14/13)^{\lfloor n/100\rfloor}$ clauses.
\end{theorem}

\begin{proof}
Let $C_1,\dots,C_S$ be a refutation, $t=\lfloor n/100\rfloor$, and call $C$ \emph{wide} if $|C^+|\ge n^2/10$. We kill the wide clauses with a random \emph{restriction}: choose a pigeon $i$ and a hole $j$ uniformly, set $p_{ij}=1$ and $p_{ij'}=0$, $p_{i'j}=0$ for all $j'\ne j$, $i'\ne i$; repeat $t$ times with the remaining pigeons and holes. The restricted variables disappear and $\mathrm{PHP}_n$ becomes $\mathrm{PHP}_{n-t}$ on the remaining $n-t+1$ pigeons and $n-t$ holes, up to clauses that become true; the restricted clauses $C_1|_\rho,\dots,C_S|_\rho$ still satisfy the hypotheses of Lemma~\ref{lem:bp} for $m=n-t$, because restricting preserves implications. A critical assignment of $\mathrm{PHP}_{n-t}$, completed by $\rho$, is a critical assignment of $\mathrm{PHP}_n$.

Fix a wide clause $C$. The disjunction $C^+$ contains at least $n^2/10$ variables $p_{ab}$. Before the $r$-th step at most $r-1\le t$ pigeons and holes have been used, and they carry at most $t(2n+1)\le\frac{n}{100}(2n+1)$ variables; so at least $n^2/10-n^2/50-n/100\ge n^2/13$ variables of $C^+$ involve only unused pigeons and holes. If the $r$-th choice $(i,j)$ is one of them, $C^+$ becomes true under $\rho$, hence so does $C$ on every critical assignment of $\mathrm{PHP}_{n-t}$: the clause is \emph{killed}, it has no pigeon in the sense of the lemma, and it stays so under further restrictions. There are at most $n(n+1)$ choices, so this happens at each step with probability at least $\frac{n^2/13}{n(n+1)}\ge\frac1{14}$ for $n\ge13$, whatever happened before. So $C$ survives all $t$ steps with probability at most $(13/14)^t$.

If $S<(14/13)^t$, then by the union bound some restriction $\rho$ kills all wide clauses. A clause that is not wide only loses variables: if $C|_\rho$ is not made true, then $(C|_\rho)^+$, computed in $\mathrm{PHP}_{n-t}$, uses only variables of $C^+$. So every clause of the restricted refutation is either killed or has $|(C|_\rho)^+|<n^2/10$. But $m=n-t\ge0.99n$ and $m^2/9\ge0.98n^2/9>n^2/10$, so Lemma~\ref{lem:bp} gives a clause with pigeons and $|C^+|\ge m^2/9$: a contradiction. Hence $S\ge(14/13)^t$.
\end{proof}

Our constants are poor (the proof gives $2^{n/10000}$ or so; Haken's argument gives $c^n$ with a better $c$), but the shape is the point: every refutation must contain a clause that talks about a constant fraction of all pigeon-hole pairs, and a random restriction would have killed such clauses if there were few of them. Eli Ben-Sasson and Avi Wigderson showed in 2001 that this is general: a short resolution refutation can always be made \emph{narrow}, using only clauses with few literals, so that lower bounds on the width give lower bounds on the size (quoted). Stronger systems are another matter. In the \emph{Frege} systems of textbook logic, where lines are arbitrary formulas, Samuel Buss found in 1987 refutations of $\mathrm{PHP}_n$ of polynomial size, by formalising counting (quoted); and no superpolynomial lower bound is known for Frege systems. The programme stops where the circuit programme stops, and for related reasons: lower bounds for a proof system tend to come from lower bounds for the circuits its lines can express.

\section{Barriers}\label{sec:barriers}

By the mid-1990s three kinds of lower bound arguments had succeeded: diagonalisation, which relativises and so cannot settle $\mathsf P$ versus $\mathsf{NP}$ (Theorem~\ref{thm:bgs}); random restrictions and polynomials for constant depth; approximation for monotone circuits. Razborov and Rudich looked at the last two and isolated what they have in common.

Identify a Boolean function $f$ of $n$ variables with its truth table, a string of length $N=2^n$. A \emph{property} is a sequence $\Gamma=(\Gamma_n)$ of sets of functions of $n$ variables.

\begin{definition}\label{def:natural}
A property $\Gamma$ is \emph{natural against} $\mathsf P/\mathsf{poly}$ if it is
\begin{itemize}
\item \emph{constructive}: whether $f\in\Gamma_n$ can be decided in time $2^{O(n)}$, polynomial in the size $N$ of the truth table;
\item \emph{large}: $\Gamma_n$ contains at least a fraction $2^{-O(n)}$ of all functions of $n$ variables;
\item \emph{useful}: for every $c$, if $n$ is large enough, no function in $\Gamma_n$ has a circuit of size $n^c$.
\end{itemize}
\end{definition}

A useful property is a lower bound proof in disguise: exhibit a function in $\mathsf{NP}$ with property $\Gamma_n$ for each $n$, and $\mathsf{NP}\not\subseteq\mathsf P/\mathsf{poly}$. The proof of Theorem~\ref{thm:rs}, for instance, shows that being far from all polynomials of degree $\sqrt n$ over $\mathbb F_3$ is useful against $\mathsf{AC}^0[3]$; this property is large (Exercise~\ref{ex:largeness}), and Razborov and Rudich checked that a closely related one, defined through the ranks of certain matrices, is constructive (quoted). Most known proofs fit the pattern.

The obstacle is cryptography. A \emph{pseudorandom function family} is a family of functions $f_s:\{0,1\}^n\to\{0,1\}$, indexed by seeds $s\in\{0,1\}^n$, such that $(s,x)\mapsto f_s(x)$ has circuits of size $n^{c_0}$ for some constant $c_0$, and which is \emph{exponentially secure}: for some $\varepsilon>0$ and all large $n$, every circuit $T$ of size $2^{n^\varepsilon}$ that may query an oracle $h$ satisfies
\[
\Bigl|\,P_s\bigl(T^{f_s}=1\bigr)-P_g\bigl(T^{g}=1\bigr)\Bigr|<2^{-n^\varepsilon},
\]
for $s$ uniform and $g$ a uniformly random function. Such families can be built from one-way functions, by Oded Goldreich, Shafi Goldwasser and Silvio Micali (1986) and Johan Håstad, Russell Impagliazzo, Leonid Levin and Michael Luby (1999), and exponentially secure ones exist if, for instance, factoring is hard enough (quoted); this is the business of Tome~C4, Chapter~8.

\begin{theorem}[Razborov, Rudich 1994]\label{thm:naturalproofs}
If an exponentially secure pseudorandom function family exists, then no property is natural against $\mathsf P/\mathsf{poly}$.
\end{theorem}

\begin{proof}
Suppose $\Gamma$ is natural: decidable in time $2^{am}$ on functions of $m$ variables, and containing a fraction at least $2^{-bm}$ of them. A natural property is a distinguisher. Let $m=\lfloor n^{\varepsilon/2}\rfloor$, and let $T$ do the following with its oracle $h$: query $h$ on the $2^m$ inputs $y0^{n-m}$, $y\in\{0,1\}^m$, which gives the truth table of $h'(y)=h(y0^{n-m})$, and output $1$ if and only if $h'\in\Gamma_m$. Simulating the decision procedure by a circuit (Proposition~\ref{prop:ppoly}), $T$ has size $2^{O(m)}$, less than $2^{n^\varepsilon}$ for large $n$.

If $h=g$ is random, $h'$ is a uniformly random function of $m$ variables, and $T$ outputs $1$ with probability at least $2^{-bm}$. If $h=f_s$, then $h'(y)=f_s(y0^{n-m})$ is computed by a circuit of size $n^{c_0}$, with $s$ and the zeros wired in; as $n\le(m+1)^{2/\varepsilon}$, this is at most $m^c$ for $c=2c_0/\varepsilon+1$ and $m$ large. Usefulness, applied with this $c$, says $h'\notin\Gamma_m$, and $T$ outputs $0$. The advantage of $T$ is at least $2^{-bm}\ge2^{-b\,n^{\varepsilon/2}}$, more than $2^{-n^\varepsilon}$ for large $n$, contradicting security.
\end{proof}

So a proof that $\mathsf{NP}$ needs large circuits must be \emph{unnatural}: it must use a property that is rare among functions, or hard to recognise. The irony, which Razborov and Rudich point out, is that the stronger our belief in cryptography, the stronger the barrier: a proof of $\mathsf P\ne\mathsf{NP}$ by the known methods would break the very pseudorandom functions that $\mathsf P\ne\mathsf{NP}$ is supposed to make possible. Williams' theorem on $\mathsf{ACC}^0$ escapes the barrier by combining an algorithm with diagonalisation, and a third barrier, \emph{algebrisation}, found by Scott Aaronson and Avi Wigderson in 2008, blocks the arithmetisation of Chapter~\ref{ch:randomness} when used alone (quoted).

The road also runs the other way. In Chapter~\ref{ch:randomness} we saw that hard functions give derandomisation. Valentine Kabanets and Russell Impagliazzo proved in 2004 the converse kind of statement: if polynomial identity testing, which has the simple randomised algorithm of Tome~C2, Chapter~8, has a deterministic polynomial algorithm, then either $\mathsf{NEXP}\not\subseteq\mathsf P/\mathsf{poly}$ or the permanent has no arithmetic circuits of polynomial size (quoted). Derandomising a simple algorithm would prove circuit lower bounds that seem far out of reach; algorithms and lower bounds are two faces of one problem.

\begin{beauty}
Shannon's counting shows that hardness is everywhere, and the whole difficulty is to put a finger on it. Each success of this chapter is a way of seeing a small circuit as something simple: a polynomial of low degree, a short list of cliques, a proof whose every line speaks of few pigeons; and each time the hard function is too rich to be simple in that sense. Then the barrier: if the property that separates simple from complex is itself simple and common, it can be fooled by pseudorandom functions. Hardness, randomness and proof turn out to be one subject, and the question of Cook and Levin is still open, sharper and better understood than ever.
\end{beauty}

\section{Exercises}

\begin{exercise}[$\star$]\label{ex:paritysmall}
Show that $\mathrm{PAR}_n$ has a circuit with $\land,\lor,\lnot$ of size at most $5(n-1)$ and a formula of size $O(n^2)$. Compare with the $2^{n-1}$ clauses of every CNF for parity.
\end{exercise}

\begin{exercise}[$\star\star$]\label{ex:everyfn}
Show that every Boolean function of $n\ge1$ variables has a circuit of size at most $5\cdot2^{n-1}$. \emph{Use} $f=(x_n\land f_1)\lor(\lnot x_n\land f_0)$.
\end{exercise}

\begin{exercise}[$\star\star$]\label{ex:majparity}
Let $\mathrm{MAJ}_m(z)=1$ when at least $m/2$ of the $m$ bits of $z$ are $1$. Show that if $\mathrm{MAJ}$ had circuits of depth $d$ and size $s(m)$ with gates $\land,\lor,\lnot,\mathrm{MOD}_3$, then $\mathrm{PAR}_n$ would have circuits of depth $d+3$ and size $O(n\,s(2n))$. Deduce that majority is not in $\mathsf{AC}^0[3]$.
\end{exercise}

\begin{exercise}[$\star$]\label{ex:orexact}
Show that every function $\{0,1\}^k\to\mathbb F_3$ is represented by exactly one multilinear polynomial, and that every polynomial over $\mathbb F_3$ equal to $x_1\lor\dots\lor x_k$ on $\{0,1\}^k$ has degree at least $k$. Why did Lemma~\ref{lem:approxpoly} need randomness?
\end{exercise}

\begin{exercise}[$\star$]\label{ex:rainbow}
For a uniformly random colouring of $n$ vertices with $k-1$ colours, compute the probability that a fixed set of $r$ vertices is rainbow, and show that it is at least $1-\binom r2/(k-1)$.
\end{exercise}

\begin{exercise}[$\star$]\label{ex:pluckerror}
Give a negative test graph and a sunflower $Z_1,Z_2,Z_3$ with core $Z$ such that $C_Z=1$ and $C_{Z_1}=C_{Z_2}=C_{Z_3}=0$ on it. Why can plucking never err on a positive test graph?
\end{exercise}

\begin{exercise}[$\star\star$]\label{ex:kannan}
Let $k\ge1$. Show that some language in $\Sigma_3^p$ has no circuits of size $n^k$. \emph{Consider the lexicographically first function of the first $\lceil(k+1)\log_2n\rceil$ input bits that has no circuit of size $n^k$.}
\end{exercise}

\begin{exercise}[$\star$]\label{ex:php32}
Write a resolution refutation of $\mathrm{PHP}_2$ (three pigeons, two holes).
\end{exercise}

\begin{exercise}[$\star$]\label{ex:critical}
Count the critical assignments of $\mathrm{PHP}_n$, and check that an $i$-critical assignment falsifies exactly one clause of $\mathrm{PHP}_n$. Show that $C$ and $C^+$ take the same value on critical assignments.
\end{exercise}

\begin{exercise}[$\star\star$]\label{ex:largeness}
Show that, as $n\to\infty$, the fraction of Boolean functions of $n$ variables that agree with some polynomial of degree at most $\sqrt n$ over $\mathbb F_3$ on at least $90\%$ of the points of $\{0,1\}^n$ tends to $0$. \emph{Use} $\sum_{j\le K}\binom Nj\le(eN/K)^K$.
\end{exercise}

\section{Hints}

\begin{hint}{ex:paritysmall}
$a\oplus b=(a\lor b)\land\lnot(a\land b)$. For the formula, balance the tree.
\end{hint}
\begin{hint}{ex:everyfn}
Count gates: $s(n)\le2s(n-1)+4$, and functions of one variable need at most one gate.
\end{hint}
\begin{hint}{ex:majparity}
Padding with constants turns majority into the threshold function ``at least $k$ ones''; then ``exactly $k$ ones'' is a threshold minus the next one.
\end{hint}
\begin{hint}{ex:orexact}
Count: there are as many functions as multilinear polynomials. What is the multilinear polynomial of $\lor$?
\end{hint}
\begin{hint}{ex:rainbow}
Colour the vertices one at a time.
\end{hint}
\begin{hint}{ex:pluckerror}
Give every petal vertex outside $Z$ the colour of a vertex of $Z$.
\end{hint}
\begin{hint}{ex:kannan}
Truth tables on $r=\lceil(k+1)\log_2n\rceil$ bits have polynomial length, so ``the circuit $C$ computes the table $T$'' is checked in polynomial time.
\end{hint}
\begin{hint}{ex:php32}
First show that pigeon $3$ is not in hole $1$.
\end{hint}
\begin{hint}{ex:critical}
Choose the missing pigeon, then a bijection.
\end{hint}
\begin{hint}{ex:largeness}
Count the polynomials, then the functions close to one polynomial.
\end{hint}

\section{Solutions}

\begin{solution}{ex:paritysmall}
\idea Parity is a chain of exclusive ors, and each exclusive or is cheap.
\writeup The identity $a\oplus b=(a\lor b)\land\lnot(a\land b)$ computes $a\oplus b$ with $4$ gates (or $5$ if one prefers a symmetric form), and $\mathrm{PAR}_n=(\dots(x_1\oplus x_2)\oplus\dots)\oplus x_n$ uses $n-1$ of them: at most $5(n-1)$ gates. A formula cannot reuse $a$ and $b$, so it doubles them: $a\oplus b=(a\land\lnot b)\lor(\lnot a\land b)$ uses each argument twice. Arrange the exclusive ors in a balanced tree of depth $\lceil\log_2n\rceil$; each leaf is then copied $2^{\lceil\log_2n\rceil}\le2n$ times, so the formula has $O(n^2)$ gates. So parity, which needs $2^{n-1}$ clauses in CNF (Tome~C1, Proposition~1.5), is easy for circuits; this is why lower bounds for depth $2$ say little about general circuits.
\end{solution}

\begin{solution}{ex:everyfn}
\idea Split on the last variable and recurse: the size at most doubles at each level.
\writeup A function of one variable is $0$, $1$, $x_1$ or $\lnot x_1$: at most one gate. For $n\ge2$ let $f_b$ be $f$ with $x_n=b$, a function of $n-1$ variables; then $f=(x_n\land f_1)\lor(\lnot x_n\land f_0)$, with $4$ gates more than circuits for $f_0$ and $f_1$. So $s(n)\le2s(n-1)+4$, that is $s(n)+4\le2\bigl(s(n-1)+4\bigr)$, and $s(n)+4\le2^{n-1}(s(1)+4)\le5\cdot2^{n-1}$. Lupanov's construction gains a factor $n$ by computing once all the functions of the last few variables and sharing them.
\end{solution}

\begin{solution}{ex:majparity}
\idea Majority with constants counts; counting gives parity.
\writeup Let $\mathrm{TH}_k(x)=1$ when at least $k$ of the $n$ bits of $x$ are $1$. For $0\le k\le n$, feed $\mathrm{MAJ}_{2n}$ with $x$, $n-k$ ones and $k$ zeros: the number of ones is at least $n$ if and only if $x$ has at least $k$ ones, so $\mathrm{TH}_k(x)=\mathrm{MAJ}_{2n}(x,1^{n-k},0^k)$, and $\mathrm{TH}_{n+1}=0$. Then ``exactly $k$ ones'' is $\mathrm{TH}_k\land\lnot\mathrm{TH}_{k+1}$, and $\mathrm{PAR}_n=\bigvee_{k\text{ odd}}(\mathrm{TH}_k\land\lnot\mathrm{TH}_{k+1})$. This adds a layer of $\lnot$, a layer of $\land$ and one $\lor$ gate: depth $d+3$, and size at most $(n+1)s(2n)+O(n)$. If majority had polynomial size circuits of constant depth with $\mathrm{MOD}_3$ gates, so would parity, contradicting Theorem~\ref{thm:rs}.
\end{solution}

\begin{solution}{ex:orexact}
\idea A polynomial that is exact on the cube is unique once multilinear, and $\lor$ needs the top monomial.
\writeup On $\{0,1\}$, $x^2=x$, so every polynomial agrees on the cube with a multilinear one of no larger degree. Every function $f$ is represented: $f=\sum_af(a)\prod_{i:a_i=1}x_i\prod_{i:a_i=0}(1-x_i)$. There are $3^{2^k}$ functions and $3^{2^k}$ multilinear polynomials (one coefficient in $\mathbb F_3$ per subset of variables), so representation is a bijection and is unique. The representative of $\lor$ is $1-\prod_i(1-x_i)$, whose coefficient of $x_1\cdots x_k$ is $\pm1\ne0$: degree $k$. A polynomial of smaller degree equal to $\lor$ on the cube would give a second, multilinear representative of degree less than $k$. So no low-degree polynomial is \emph{exact} for $\lor$ gates of large fan-in, and the lemma settles for polynomials that are right at most points, which a random choice provides.
\end{solution}

\begin{solution}{ex:rainbow}
\idea Colour the vertices one after the other and avoid the colours already used.
\writeup The $i$-th vertex of the set, $i=1,\dots,r$, must avoid the $i-1$ colours already used, so the probability is $\prod_{i=0}^{r-1}\bigl(1-\frac i{k-1}\bigr)$ (zero if $r>k-1$). Alternatively, the set fails to be rainbow only if one of the $\binom r2$ pairs has equal colours, each with probability $\frac1{k-1}$; the union bound gives at least $1-\binom r2/(k-1)$. With $r=\ell=\sqrt k$ this is at least $\frac12$, as used in the proof of Theorem~\ref{thm:monotoneclique}.
\end{solution}

\begin{solution}{ex:pluckerror}
\idea A rainbow core does not make the petals rainbow.
\writeup Let $k-1=3$, $Z=\{1,2\}$ and $Z_i=\{1,2,2+i\}$ for $i=1,2,3$, a sunflower with core $Z$. Colour $1$ red, $2$ blue, and $3,4,5$ red. Then $Z$ is rainbow, so $C_Z=1$, while each $Z_i$ contains two red vertices, so $C_{Z_i}=0$. On a positive test graph with clique $K$, $C_{Z_i}=1$ means $Z_i\subseteq K$, which implies $Z\subseteq K$; so replacing the petals by the core can only change values from $0$ to $1$, and the circuit's value there is $1$ anyway whenever it matters.
\end{solution}

\begin{solution}{ex:kannan}
\idea Counting gives a hard function on few bits, and three quantifiers find the first one.
\writeup Let $r=\lceil(k+1)\log_2n\rceil$, so $2^r\ge n^{k+1}$. By the counting of Theorem~\ref{thm:shannon}, there are at most $2^{O(n^k\log n)}$ circuits of size $n^k$ with $r$ inputs, fewer than the $2^{2^r}$ functions of $r$ bits for large $n$; so a function of $r$ bits with no such circuit exists, and we let $T_n$ be the lexicographically first one, as a truth table of length $2^r$, polynomial in $n$. Let $L=\{x:T_{|x|}(x_1\dots x_r)=1\}$. A circuit of size $n^k$ for $L\cap\{0,1\}^n$, with $x_{r+1},\dots,x_n$ fixed to $0$, would compute $T_n$: so $L$ has no circuits of size $n^k$. And $x\in L$ if and only if there is a table $T$ such that $T(x_1\dots x_r)=1$, every circuit $C$ of size $n^k$ differs from $T$, and every table $T'<T$ is computed by some circuit $C'$ of size $n^k$. Comparing a circuit with a table takes polynomial time, since the table has polynomial length. The condition reads $\exists T\,\forall C\,\forall T'\,\exists C'$, with polynomial-time matrix, which is $\Sigma_3^p$ (the two universal quantifiers merge). Ravi Kannan proved in 1982 the sharper statement for $\Sigma_2^p$, combining this argument with Theorem~\ref{thm:karplipton}.
\end{solution}

\begin{solution}{ex:php32}
\idea Follow the human argument: pigeon $3$ fits in neither hole.
\writeup Write $P_i=p_{i1}\lor p_{i2}$. Resolving $\lnot p_{31}\lor\lnot p_{11}$ with $P_1$ gives $\lnot p_{31}\lor p_{12}$; resolving $\lnot p_{31}\lor\lnot p_{21}$ with $P_2$ gives $\lnot p_{31}\lor p_{22}$; resolving the first with $\lnot p_{12}\lor\lnot p_{22}$ gives $\lnot p_{31}\lor\lnot p_{22}$, and with the second, $\lnot p_{31}$. Symmetrically, with the hole clauses of hole $2$ and then hole $1$, derive $\lnot p_{32}$. Resolving $P_3=p_{31}\lor p_{32}$ with $\lnot p_{31}$ gives $p_{32}$, and with $\lnot p_{32}$ the empty clause: ten derived clauses in all.
\end{solution}

\begin{solution}{ex:critical}
\idea A critical assignment is a perfect placement with one pigeon left out.
\writeup Choose the missing pigeon ($n+1$ ways) and a bijection from the other pigeons to the holes ($n!$ ways): $(n+1)!$ critical assignments. Each hole contains exactly one pigeon, so all hole clauses hold, and every pigeon but $i$ sits somewhere, so the only false clause is the pigeon clause of $i$. On a critical assignment, $\lnot p_{ij}$ holds if and only if hole $j$, which is occupied, is occupied by a pigeon $i'\ne i$, that is if and only if $\bigvee_{i'\ne i}p_{i'j}$ holds; so each literal of $C$ and its replacement in $C^+$ have the same value, and so do $C$ and $C^+$.
\end{solution}

\begin{solution}{ex:largeness}
\idea There are few low-degree polynomials and few functions near each of them.
\writeup Let $N=2^n$. As in Exercise~\ref{ex:orexact}, on the cube a polynomial of degree at most $\sqrt n$ agrees with a multilinear one of degree at most $\sqrt n$, and there are $3^B$ of these, with $B=\sum_{i\le\sqrt n}\binom ni\le(n+1)^{\sqrt n}$. A function that agrees with a fixed polynomial on at least $90\%$ of the points differs from it on a set of at most $K=N/10$ points, on which its value is then determined; so there are at most $\sum_{j\le K}\binom Nj\le(10e)^{N/10}\le2^{0.48N}$ such functions. The fraction of all $2^N$ functions that are close to some low-degree polynomial is at most $3^B2^{0.48N}2^{-N}=2^{B\log_23-0.52N}$, which tends to $0$ since $B\le2^{\sqrt n\log_2(n+1)}$ is tiny compared with $N=2^n$. So the property ``far from all polynomials of degree $\sqrt n$'' holds for almost all functions: it is large.
\end{solution}

\section*{Sources and further reading}

\begin{itemize}
\item \emph{A research paper to read.} A.~A.~Razborov and S.~Rudich, \emph{Natural proofs}, Journal of Computer and System Sciences 55 (1997), 24--35 (STOC 1994). Twelve pages, and one of the few papers about why a problem is hard to solve rather than about a solution. Read the introduction for the definitions of constructivity, largeness and usefulness, then the section that checks, one after the other, that the known lower bounds (for $\mathsf{AC}^0$, for $\mathsf{AC}^0[p]$, for monotone circuits) are natural: compare it with our Sections~\ref{sec:parity} and~\ref{sec:monotone}. The main theorem is our Theorem~\ref{thm:naturalproofs}, with a sharper accounting of the parameters. The final section discusses ways out; it is worth rereading after Williams' theorem, which takes one of them. The paper received the G\"odel Prize in 2007.
\item \emph{The stories.} C.~E.~Shannon, \emph{A symbolic analysis of relay and switching circuits}, Transactions of the American Institute of Electrical Engineers 57 (1938), 713--723; C.~E.~Shannon, \emph{The synthesis of two-terminal switching circuits}, Bell System Technical Journal 28 (1949), 59--98; M.~Sipser, \emph{The history and status of the P versus NP question}, Proceedings of the 24th ACM Symposium on Theory of Computing (1992), 603--618.
\item O.~B.~Lupanov, \emph{A method of circuit synthesis}, Izvestiya VUZ Radiofizika 1 (1958), 120--140; J.~Li and T.~Yang, \emph{$3.1n-o(n)$ circuit lower bounds for explicit functions}, Proceedings of the 54th ACM Symposium on Theory of Computing (2022), 1180--1193; R.~M.~Karp and R.~J.~Lipton, \emph{Some connections between nonuniform and uniform complexity classes}, Proceedings of the 12th ACM Symposium on Theory of Computing (1980), 302--309; R.~Kannan, \emph{Circuit-size lower bounds and non-reducibility to sparse sets}, Information and Control 55 (1982), 40--56.
\item M.~Furst, J.~B.~Saxe and M.~Sipser, \emph{Parity, circuits, and the polynomial-time hierarchy}, Mathematical Systems Theory 17 (1984), 13--27; M.~Ajtai, \emph{$\Sigma^1_1$-formulae on finite structures}, Annals of Pure and Applied Logic 24 (1983), 1--48; J.~Håstad, \emph{Almost optimal lower bounds for small depth circuits}, Proceedings of the 18th ACM Symposium on Theory of Computing (1986), 6--20; A.~A.~Razborov, \emph{Lower bounds on the size of bounded depth circuits over a complete basis with logical addition}, Mathematical Notes 41 (1987), 333--338; R.~Smolensky, \emph{Algebraic methods in the theory of lower bounds for Boolean circuit complexity}, Proceedings of the 19th ACM Symposium on Theory of Computing (1987), 77--82; R.~Williams, \emph{Nonuniform ACC circuit lower bounds}, Journal of the ACM 61 (2014), article 2.
\item A.~A.~Razborov, \emph{Lower bounds on the monotone complexity of some Boolean functions}, Soviet Mathematics Doklady 31 (1985), 354--357; A.~A.~Razborov, \emph{Lower bounds on monotone complexity of the logical permanent}, Mathematical Notes 37 (1985), 485--493; N.~Alon and R.~B.~Boppana, \emph{The monotone circuit complexity of Boolean functions}, Combinatorica 7 (1987), 1--22; \'E.~Tardos, \emph{The gap between monotone and non-monotone circuit complexity is exponential}, Combinatorica 8 (1988), 141--142.
\item S.~A.~Cook and R.~A.~Reckhow, \emph{The relative efficiency of propositional proof systems}, Journal of Symbolic Logic 44 (1979), 36--50; A.~Haken, \emph{The intractability of resolution}, Theoretical Computer Science 39 (1985), 297--308; P.~Beame and T.~Pitassi, \emph{Simplified and improved resolution lower bounds}, Proceedings of the 37th IEEE Symposium on Foundations of Computer Science (1996), 274--282; E.~Ben-Sasson and A.~Wigderson, \emph{Short proofs are narrow --- resolution made simple}, Journal of the ACM 48 (2001), 149--169; S.~R.~Buss, \emph{Polynomial size proofs of the propositional pigeonhole principle}, Journal of Symbolic Logic 52 (1987), 916--927.
\item O.~Goldreich, S.~Goldwasser and S.~Micali, \emph{How to construct random functions}, Journal of the ACM 33 (1986), 792--807; J.~Håstad, R.~Impagliazzo, L.~A.~Levin and M.~Luby, \emph{A pseudorandom generator from any one-way function}, SIAM Journal on Computing 28 (1999), 1364--1396; S.~Aaronson and A.~Wigderson, \emph{Algebrization: a new barrier in complexity theory}, ACM Transactions on Computation Theory 1 (2009), article 2; V.~Kabanets and R.~Impagliazzo, \emph{Derandomizing polynomial identity tests means proving circuit lower bounds}, Computational Complexity 13 (2004), 1--46.
\item S.~Jukna, \emph{Boolean Function Complexity: Advances and Frontiers}, Springer, 2012; S.~Arora and B.~Barak, \emph{Computational Complexity: A Modern Approach}, Cambridge University Press, 2009, Chapters~6, 14, 15 and~23.
\end{itemize}
